% ============================================================
%  Modification and Validation of pyBuildingEnergy for
%  Australian Residential Building Energy Simulation
%  Compile: pdflatex → pdflatex → bibtex → pdflatex → pdflatex
% ============================================================
\documentclass[12pt, a4paper]{article}

% ---- Packages -----------------------------------------------
\usepackage[utf8]{inputenc}
\usepackage[T1]{fontenc}
\usepackage[a4paper, top=2.5cm, bottom=2.5cm,
            left=3.0cm, right=2.5cm]{geometry}
\usepackage{color}
\usepackage{soul}
\usepackage{amsmath, amssymb}
\usepackage{graphicx}
\usepackage{booktabs}
\usepackage{tabularx}
\usepackage{longtable}
\usepackage{array}
\usepackage{multirow}
\usepackage[hidelinks]{hyperref}
\usepackage{xcolor}
\usepackage{listings}
\usepackage{enumitem}
\usepackage{setspace}
\usepackage{fancyhdr}
\usepackage{cite}
\usepackage{caption}
\usepackage{subcaption}
\usepackage{microtype}
\usepackage{lmodern}
\usepackage{parskip}
\usepackage{titlesec}
\usepackage{natbib}
% Simple DOI display (no external package needed)
\newcommand{\doi}[1]{\href{https://doi.org/#1}{\texttt{doi:#1}}}

% ---- Section title formatting (journal style) ---------------
\titleformat{\section}{\large\bfseries}{\thesection.}{0.5em}{}
\titleformat{\subsection}{\normalsize\bfseries}{\thesubsection.}{0.5em}{}
\titleformat{\subsubsection}{\normalsize\itshape}{\thesubsubsection.}{0.5em}{}

% ---- Spacing -----------------------------------------------
\setstretch{1.15}
\setlength{\parskip}{0.5em}

% ---- Code listing style ------------------------------------
\definecolor{codegreen}{rgb}{0,0.55,0}
\definecolor{codegray}{rgb}{0.5,0.5,0.5}
\definecolor{codepurple}{rgb}{0.5,0,0.7}
\definecolor{backcolour}{rgb}{0.97,0.97,0.97}
\definecolor{framecolour}{rgb}{0.8,0.8,0.8}

\lstdefinestyle{pythonstyle}{
    backgroundcolor=\color{backcolour},
    commentstyle=\color{codegreen}\itshape,
    keywordstyle=\color{blue!80!black}\bfseries,
    numberstyle=\tiny\color{codegray},
    stringstyle=\color{codepurple},
    basicstyle=\ttfamily\footnotesize,
    breakatwhitespace=false,
    breaklines=true,
    captionpos=b,
    keepspaces=true,
    numbers=left,
    numbersep=6pt,
    showspaces=false,
    showstringspaces=false,
    showtabs=false,
    tabsize=2,
    language=Python,
    frame=single,
    rulecolor=\color{framecolour},
    xleftmargin=12pt,
    framexleftmargin=12pt
}
\lstset{style=pythonstyle}

% ---- Page style --------------------------------------------
\pagestyle{fancy}
\fancyhf{}
\rhead{\small\itshape pyBuildingEnergy Modifications for Australian Simulation}
\lhead{\small\itshape AIB Research Team}
\cfoot{\thepage}
\renewcommand{\headrulewidth}{0.4pt}

% ---- Hyperref setup ----------------------------------------
\hypersetup{
    colorlinks=true,
    linkcolor=blue!60!black,
    citecolor=green!50!black,
    urlcolor=blue!60!black
}

% ---- Title --------------------------------------------------
\title{
    \vspace{-0.5cm}
    \LARGE\textbf{Revising pyBuildingEnergy
    Simulation Engine- A methodology for Australian Residential Buildings}
    \vspace{0.3cm}
}

\author{
    AIB Research Team\\[0.3em]
%    \small\textit{Repository: \texttt{samiraghafarigousheh-sys/PyBuildingEnergy\_AIBteam\_AU}}
}

\date{\today}

% =============================================================
\begin{document}
% =============================================================

\maketitle
\thispagestyle{fancy}

% ---- Abstract -----------------------------------------------
\begin{abstract}
\noindent
This paper documents the modifications made to pybuildingenergy as the baseline for measuring energy consumption based on the standard.... to adapt the engine for Australian residential and localized weather conditions contexts while comparing the measurement with the ISo 52016 and the more advanced energy calculation engines such as energy plus .





\end{abstract}

\vspace{0.3cm}
\noindent\textbf{Keywords:} ISO 52016-1, building energy simulation, pyBuildingEnergy,
EnergyPlus validation, Australian residential buildings

\tableofcontents
\newpage

% =============================================================
\section{Introduction}
\label{sec:introduction}
% =============================================================
 
Accurate prediction of building energy consumption is central to the design of
low-carbon residential stock and to the assessment of energy efficiency policy.
At the national level, the Australian National Construction Code (NCC) mandates
minimum performance levels for new dwellings, and assessors increasingly require
dynamic simulation evidence that these levels are met under local climatic conditions.
At the project level, energy modelling informs decisions about building fabric,
glazing, mechanical ventilation, and equipment sizing that collectively determine
whether a building is comfortable and energy-efficient throughout its life.
 
Apart from data-driven approach,Two broad classes of tools address these requirements.
Detailed dynamic simulation engines such as EnergyPlus~\cite{energyplus} and
TRNSYS~\cite{trnsys} solve the full heat and mass balance of a building at
sub-hourly resolution, accounting for variable surface heat transfer coefficients,
hourly window optical properties, and multi-zone airflow.
Their fidelity is high, but they require extensive parameterisation, specialist
expertise, and significant computational time.
Simplified hourly methods, codified in international standards, offer a middle
ground: they are computationally lightweight and can be run on commodity hardware
with modest user effort, making them suitable for compliance screening, early design
exploration, and large-scale stock analyses.
 
The international Organization for Standardization (ISO)~52016-1:2018 standard~\cite{iso52016} is the current European and
internationally adopted simplified hourly method, replacing the earlier EN~ISO~13790
5R1C single-capacitance model with a multi-node resistance--capacitance network that
more accurately represents the thermal mass and dynamic behaviour of building
envelopes.
The open-source Python library \texttt{pyBuildingEnergy}~\cite{eurac2024},
developed by EURAC Research, provides an executable implementation of this standard
and is the simulation backbone adopted in the present study.
 
The \texttt{pyBuildingEnergy} library was designed primarily for
European building stock and regulatory contexts.  Despite its advantages, there exists some limitations in the context of Australian
residential apartments. \hl{As per \textit{author et al.,} method-level simplifications in ISO~52016-1 can produce systematic cooling inaccuracy (citation). Additionally, it is found that the library 
lacks Australia -specific features such as ....Also, exsiting implementation defects in the engine cause
systematic errors regardless of the standard's method-level accuracy.}
 
The objective of this paper is therefore to develop a more sutable and reproducible methodology to address these limitations.
The methodology is anchored to a specific case study building in Melbourne, Australia — and validated against EnergyPlus as the reference benchmark.

The result is a modified implementation, released as \texttt{AIB}, whose
departures from the upstream release are individually documented, traceable to
the provision of the Standard or the measurement from which they derive, and
subject to an energy-balance closure criterion applied uniformly across the
correction sequence.
 
\section{Literature Review}
\subsection{From ISO~13790 to ISO~52016-1}
ISO~52016-1 was developed under the European Committee for Standardization (CEN) Mandate M/480, in support of the recast Energy
Performance of Buildings Directive (2010/31/EU), as the successor to the withdrawn
EN~ISO~13790 simple hourly method.
While ISO~13790 represents building as a single 5-resistance/1-capacitance (5R1C)
network, ISO~52016-1 discretises every building element into a multi-node
resistance--capacitance chain and resolves the zone air-node energy balance at every
hourly timestep.
This gives a more detailed representation of envelope thermal mass and dynamic
behaviour than its predecessor, while remaining far less data-intensive than a
fully dynamic simulation engine.
 

As a reference, a growing body of literature uses these fully dynamic simulation engines to validate and find the specific source for the discrepancy between the results.
For example, \cite{zakula2021} compared the Standard against TRNSYS in 147,456 cases
that span eight climate zones developed by the American Society of Heating, Refrigerating and Air-Conditioning Engineers (ASHRAE), and identified use of constant,
angle-independent window optical properties as the leading cause of cooling
inaccuracy, specifically for ASHRAE zone 6, cold Nordic climate.
Recomputing the window's total solar energy transmittance $g_{\text{win}}$ hourly
as a function of solar angle of incidence raised the share of acceptable
cooling cases to 35\% for Zone~6 and to 75--100\% for warmer moderate climates.
 
\cite{ballarini2020} and \cite{ballarini2019} reached a complementary conclusion
using EnergyPlus as the reference, attributing residual cooling error to the
Standard's use of constant convective and radiative surface heat transfer
coefficients inherited from ISO~13789, and to differences in the modelling of
long-wave radiative exchange with the sky.
\cite{deluca2022} isolated the effect of each correction introduced by the Italian
National Annex to ISO~52016-1 and found that solar-angle- and time-dependent
correction to the window's solar energy transmittance produced by far the largest
change in predicted heating and cooling needs, corroborating \cite{zakula2021}'s
finding.
\cite{summa2022} compared the Standard's own conduction discretisation (Annex~B)
against the Italian Annex~A alternative and likewise found lower error for heating
than for cooling, with RMSE degrading as glazing solar transmittance is reduced.

\cite{campana2019} benchmarked ISO~52016-1 against BESTEST/ASHRAE~140 reference
cases and found that for heavyweight constructions, the Standard's five-class
thermal-mass scheme produced deviations from the BESTEST range exceeding 90\%,
attributing this to the use of discrete capacitance classes rather than
construction-specific areal heat capacities.

Across all these studies, three recurring method-level causes of error emerge:
(1)~constant, hourly-invariant surface heat transfer coefficients; (2)~constant,
angle-independent window optical properties; and (3)~a classified rather than
physically calculated representation of envelope thermal mass.
 
 
Overall, the existing studies characterise the accuracy of the Standard on the assumption
of a correct implementation.
However, limited attempts have been made to address the additional layer of defects that can be introduced when the Standard is translated into an open-source tool.

Modifying the regional features, which is the main objective of this paper, can only be applied and assessed meaningfully after the underlying calculation has been checked. 
A difference between two programs may arise from an implementation error, a simplified physical method or inconsistent inputs. Each cause requires a different response. For this reason, the study examines individual heat-transfer terms as well as total energy demand.

The investigation proceeds in three stages. First, we identify and repair implementation errors and inconsistent interpretations of the inputs. Second, we test two modifications motivated by the ISO 52016 literature: angular solar transmission and wind-dependent exterior convection. These changes are intended to improve the physical representation and prepare the model for Australian applications, but their accuracy must be demonstrated rather than assumed. Third, we examine Australian specific construct and weather features including leakage assumptions, Melbourne wind exposure and alternative construction-layer descriptions.

% =============================================================
\section{Methodology}
\label{sec:methodology}
\label{sec:case_study}
% =============================================================
\subsection{Case study building}
\label{sec:building}

The case study is Apartment 305 at 50 Barry Street, Carlton, Melbourne, recorded as constructed in 2006. It is a 20 m² studio with a height of 2.7 m and an air volume of 54 m³. Its west-facing external facade contains 11.88 m² of opaque wall and two windows with a combined area of 1.62 m². The remaining walls, floor and ceiling adjoin other parts of the building. Figures~\ref{fig 1} and~\ref{fig 2} show the case documentation and model geometry, and Table~\ref{tab:building} lists the principal inputs.


 
\begin{figure}
    \centering
    \includegraphics[width=1\linewidth]{Figures/305, 50 Barry Street.png}
    \caption{Case-study photographs, plan and thermal-boundary.}
    \label{fig 1}
\end{figure}




\begin{figure}
    \centering
    \includegraphics[width=1\linewidth]{Figures/F13_case_geometry.png}
    \caption{case study geometry schematic reproduced from the supplied case documentation.}
    \label{fig 2}
\end{figure}


\begin{table}[htbp]
\centering
\caption{Principal inputs for the diagnostic apartment}
\label{tab:building}
\begin{tabular}{ll}
\toprule
Parameter & Value \\
\midrule
Location & latitude $-37.800$, longitude $144.968$ \\
Net floor area & 20\,m$^2$ \\
Ceiling height & 2.7\,m \\
Exterior wall $U$-value & 1.00\,W/(m$^2$\,K), brick veneer, R1.0 \\
Interior partition $U$-value & 2.50\,W/(m$^2$\,K) \\
Interior slab $U$-value & 1.80\,W/(m$^2$\,K) \\
Window $U$-value / $g_{gl,n}$ & 5.40\,W/(m$^2$\,K) / 0.65 \\
Window orientation & west (270\textdegree), 1.62\,m$^2$ \\
External solar absorptance & 0.75 \\
Adjacent zones & 5, conditioned at 20\,\textdegree C \\
Exposed perimeter & 0\,m \\
Setpoints & 18\,/\,26\,\textdegree C (setback 15\,/\,28) \\
Ventilation & 2.0\,l/(s\,m$^2$) \\
HVAC & ideal loads (energy need, not system) \\
\bottomrule
\end{tabular}
\end{table}

Heat transfer to neighbouring spaces is particularly important in this apartment. The five shared opaque boundaries total 75.1 m², or about 84.8\% of the envelope area. 
With the diagnostic U-values, these boundaries account for about 88.6\% of the sum of U-value multiplied by area. 

The primary diagnostic model also assigns zero heat capacity to opaque elements. This deliberately removes heat storage so that implementation and boundary-condition effects can be examined more directly.
%%%%%%%%%%%%%%%%%%%%%%%%%%%%%%%%%%%%%%%%%% 
\subsection{Weather data}
\label{sec:weather}
%%%%%%%%%%%%%%%%%%%%%%%%%%%%%%%%%%%%%%%%%%%%%%
Both engines use the same hourly weather record from Melbourne Essendon Fields. We selected this file after checking the initially considered Melbourne Regional Office record. The checks covered the number of records, missing-value codes and monthly weather patterns. In the archived Regional Office file, 2941 of 2952 hours across January, March, July and September have zero wind speed. These nearly calm monthly blocks raised a data-quality concern even though the annual file was complete.

The Essendon file has an annual mean station wind speed of 4.84 m/s and does not show the same near-zero monthly blocks. This matters because a wind-dependent calculation uses those values directly to determine heat exchange.

For exterior convection, we adjust station wind to a facade height of 6.75 m using a profile that accounts for the terrain at the station and at the building. The adopted suburban site assumptions give a multiplier of 0.6574 and an annual mean local wind speed of about 3.18 m/s. Wind direction, shelter from nearby buildings and facade pressure coefficients are not resolved.


%%%%%%%%%%%%%%%%%%%%%%%%%%%%%%%%%%%%%%%%%%%%%%%%%%%%%%%%%%%%
\subsection{Simulation engines and diagnosing discrepancies}
\label{sec:engines}
%%%%%%%%%%%%%%%%%%%%%%%%%%%%%%%%%%%%%%%%%%%%%%%%%%%%%%%%%%%%%%
Two engines are used, driven from a single shared building definition:
 
\begin{enumerate}
\item \textbf{ISO 52016-1} as implemented in \texttt{pyBuildingEnergy}, in two
  configurations --- the unmodified release, which serves as the baseline, and
  the modified engine developed gradually in 3 stages mentioned in the methodology.
\item \textbf{EnergyPlus 24.1.0}, driven by a generated Input Data File and
  adopted as the high-fidelity reference benchmark.
\end{enumerate}

Both engines are run with ideal loads, indicating that the reported quantity is the
building’s pure energy need.

Before comparing results, we checked that both models represented the same apartment and operating conditions. The checks covered geometry, window and wall areas, orientation, weather, schedules and neighbouring temperatures. 

We then traced the main heat-transfer calculations from the inputs, through the code, to the reported outputs. For airflow, this involved checking whether ventilation and infiltration were both included, whether incoming air was assigned the correct temperature, and which surface areas were used to estimate outdoor leakage. Similar checks examined internal gains, ground contact and neighbouring spaces which are the basis of implementation repaire in the next section.

To support this review, we reconstructed a heat balance from its individual terms. The balance accounts for heat entering or leaving the zone air and the represented fabric, together with changes in stored energy.

The residual of the heat balance was defined as
\begin{equation}
  r = \Phi_{\mathrm{HVAC}} + \Phi_{\mathrm{int}} + \Phi_{\mathrm{sol}}
      + \Phi_{a} + \Phi_{\mathrm{bnd}} - \frac{\Delta U}{\Delta t},
  \label{eq:residual}
\end{equation}
where $\Phi_{\mathrm{HVAC}}$ is the sensible heating or cooling power,
$\Phi_{\mathrm{int}}$ and $\Phi_{\mathrm{sol}}$ are the internal and solar
gains, $\Phi_{a}$ is the heat carried by incoming and outgoing air, and
$\Phi_{\mathrm{bnd}}$ is the heat transfer across external, adjacent and
ground boundaries. The term $\Delta U/\Delta t$ represents the change in
stored heat per unit time. The residual $r$ is the difference left after
these terms are combined. A residual close to zero shows that the reported
heat flows balance. However, a calculation can balance while using an
incorrect area or source temperature. We therefore checked those inputs and
equations separately rather than relying on balance closure alone.

Finally, controlled runs examined selected mechanisms, such as infiltration,
while keeping the remaining settings unchanged. Disabling and enabling
infiltration, for example, allowed us to check its contribution to airflow
conductance and energy demand. Earlier sweeps of ventilation rate and fixed
airflow conductance were used to explore sensitivity. A change in demand
during such a sweep does not by itself identify a software error. Our
diagnosis therefore combines inspection of the calculation with controlled
output comparisons, rather than adjusting inputs solely to reproduce the
EnergyPlus annual total. 

%%%%%%%%%%%%%%%%%%%%%%%%%%%%%%%%%%%%%%%%%%%%%%%%%%%
\subsection{Stage 1:engine problems}
\label{sec:engine problems}
%%%%%%%%%%%%%%%%%%%%%%%%%%%%%%%%%%%%%%%%%%%%%%%%%%% 
The first stage examines whether the original pyBuildingEnergy implementation uses the building inputs as intended. This is necessary because a difference in annual demand may arise from incorrect input handling or missing model capabilities, as well as from differences between simulation methods. 

The checks compare the specified inputs with the quantities used during simulation. They address internal heat gains, neighbouring-space conditions, ground contact, exposed-envelope area, ventilation and infiltration. Hourly heat flows and stored temperature states are also examined to identify inconsistencies in energy accounting.

Verified corrections are then retained in a corrected version of pyBuildingEnergy. Appendix A documents each change and distinguishes implementation repairs from added model capabilities.

The original and corrected versions are evaluated against the same primary EnergyPlus reference. This comparison establishes the combined effect of correcting the implementation and supporting the intended building specification.
The corrected version provides the starting point for the next two stages. Therefore, a subsequent variant with the solar and convection modifications disabled still contains the implementation corrections.



%%%%%%%%%%%%%%%%%%%%%%%%%%%%%%%%%%%%%%%%%%%%%%%%%%%
\subsection{Stage 2: general ISO limitations from literature}
\label{sec:general ISO limitations from literature}
%%%%%%%%%%%%%%%%%%%%%%%%%%%%%%%%%%%%%%%%%%%%%%%%%%% 

The literature reviewed in Section~2 identifies three recurring method-level
simplifications that can contribute to differences between ISO~52016-1 and
detailed dynamic simulation: the use of fixed surface heat-transfer
coefficients, angle-independent glazing properties, and a simplified
classification of envelope thermal mass.

In the present study, two of these limitations are evaluated directly:
the treatment of solar transmittance through glazing and the use of fixed
external convective heat-transfer coefficients. Both affect the interaction
between the building envelope and hourly outdoor conditions and can therefore
be examined directly against the same EnergyPlus reference model.

The third limitation identified in the literature concerns the simplified
classification of envelope thermal mass, particularly for heavyweight
constructions. This modification is not adopted directly in the present study.
The target Australian residential stock is predominantly lightweight, while
common brick-veneer construction places much of its masonry mass outside the
insulated part of the envelope, reducing its contribution to indoor thermal
storage. Refinement of the heavyweight thermal-mass classification was therefore
considered less relevant to the present application.

Construction representation is nevertheless important when adapting the model
to Australian residential buildings. Rather than modifying the
ISO~52016-1 thermal-mass formulation itself, this issue is addressed separately
in Stage~3 through Australian-specific construction-layer descriptions. The two
approaches therefore address different aspects of the problem: the
literature-based thermal-mass modification concerns how thermal capacity is
represented mathematically within the ISO calculation, whereas the Stage~3
refinement concerns the physical construction supplied to that calculation.

\subsubsection{Angular solar transmittance}

The first modification concerns solar radiation transmitted through glazing.
In the fixed treatment, the fraction of incident solar radiation transmitted
through a window is represented by a constant multiplier. However, glazing
transmittance varies with the angle at which solar radiation reaches the
surface.

The modified calculation therefore evaluates the angular dependence of direct
solar radiation at each hourly timestep. Direct radiation is corrected according
to the solar incidence angle, while diffuse and ground-reflected components are
treated separately using an appropriate diffuse-radiation factor. This replaces
the fixed solar multiplier used in the comparison model while retaining the
same glazing properties and window geometry.

The purpose of this modification is not to assume that an angle-dependent
treatment will necessarily improve agreement with EnergyPlus, but to isolate
the effect of the approximation identified in previous comparisons of
ISO~52016-1 with detailed simulation methods
\cite{zakula2021,deluca2022}.

This issue is particularly relevant to the present application because solar
gain through glazing can contribute substantially to cooling demand. The
case-study apartment has a west-facing glazed facade, for which the incidence
angle changes considerably over the day. An hourly treatment therefore allows
the model to respond to the actual solar geometry rather than applying the same
transmission factor throughout the year.

\subsubsection{Wind-dependent external convection}

The second modification concerns convective heat exchange between
air-exposed surfaces and the outdoor environment. The fixed formulation uses
a constant external convective heat-transfer coefficient. In reality, this
exchange varies with wind speed, so a constant coefficient cannot represent
hour-to-hour changes in the external surface boundary condition.

The modified calculation therefore evaluates the external convective
coefficient as a function of hourly local wind speed for exposed walls and
windows. The wind speed measured in the weather file is first adjusted from
the meteorological station conditions to the height and terrain of the
case-study facade, as described in Section~\ref{sec:weather}. The resulting
local wind speed is then used in the hourly external convection calculation.

The modification is applied only to the external surface boundary condition.
The standard coefficient used when deriving the fixed construction resistance
from the rated $U$-value is retained. This distinction is necessary because
construction resistance is a material property and should not change with
weather, whereas convective exchange at the exterior surface does vary with
wind.

Only surfaces exposed to outdoor air are assigned the wind-dependent
coefficient. Adjacent-zone and ground-coupled surfaces retain their existing
boundary treatments.

This modification is relevant to the Australian application because wind
conditions can vary substantially over time and location. In Melbourne, the
difference between station wind speed and the wind actually experienced at the
building facade is also important. For this reason, the effect of the
wind-dependent formulation is evaluated using the terrain- and height-adjusted
local wind rather than applying the meteorological station value directly.

\subsubsection{Evaluation of the two modifications}

To separate the effects of solar transmittance and external convection,
four \texttt{pyBuildingEnergy} variants are evaluated
(Table~\ref{tab:variants}). All variants retain the implementation corrections
identified in Stage~1 and use the same building geometry, schedules, weather
data, ventilation assumptions and boundary conditions. Each variant is compared
against the same unchanged EnergyPlus reference.

\begin{table}[htbp]
\centering
\caption{Variants used to evaluate the two literature-based physical modelling
modifications.}
\label{tab:variants}
\begin{tabular}{p{3.6cm}p{4.6cm}p{4.6cm}}
\toprule
\textbf{Variant} &
\textbf{Solar transmittance treatment} &
\textbf{External convection treatment} \\
\midrule
Corrected fixed methods &
Fixed solar multiplier &
Fixed coefficient \\

Angular transmittance only &
Dependent on solar incidence angle &
Fixed coefficient \\

Dynamic convection only &
Fixed solar multiplier &
Dependent on local wind speed \\

Both modifications &
Dependent on solar incidence angle &
Dependent on local wind speed \\
\bottomrule
\end{tabular}
\end{table}

The corrected fixed-method variant provides the reference state for this stage.
The two single-modification variants show the independent effect of each
change, while the final variant shows their combined effect. This arrangement
allows the contribution of each modification to heating and cooling energy
need to be identified without assuming in advance that either modification
will improve agreement with EnergyPlus.

The results of these comparisons are reported separately in the Results
section, including the effect of each modification when applied alone and
their combined effect. Appendix~B provides the governing equations and
implementation details, including the angular glazing formulation,
beam--diffuse separation, external convective correlation, and conversion from
meteorological station wind speed to local facade wind speed.

%%%%%%%%%%%%%%%%%%%%%%%%%%%%%%%%%%%%%%%%%%%%%%%%%%%
\subsection{Stage 3: Australian-specific refinements}
\label{sec:AU_refinements}
%%%%%%%%%%%%%%%%%%%%%%%%%%%%%%%%%%%%%%%%%%%%%%%%%%%

The first two stages address the reliability of the calculation itself:
Stage~1 verifies and corrects the implementation, while Stage~2 evaluates
selected method-level simplifications identified in the literature. The third
stage addresses a different question: whether the inputs and assumptions used
by the corrected model are representative of Australian residential buildings.

Two regional refinements are considered. The first concerns envelope air
leakage, for which the default values inherited from a European-oriented
building stock may not represent Australian dwellings. The second concerns
the physical description of opaque constructions, for which representative
Australian material-layer assemblies are introduced instead of relying only
on generic construction descriptions.

These refinements are not treated as corrections to ISO~52016-1 itself.
Rather, they adapt the information supplied to the calculation to the
Australian context. Their effects are evaluated after the Stage~1 corrections
have been retained and the Stage~2 physical modifications have been established.

% ---------------------------------------------------
\subsubsection{Australian envelope-leakage parameterisation}
% ---------------------------------------------------

The first refinement concerns the permeability of the building envelope.
Air leakage is strongly influenced by construction practice, building age
and airtightness requirements. These characteristics vary between national
building stocks, so default leakage values developed for one context cannot
necessarily be transferred directly to another.

This issue is particularly important for Australian residential modelling
because measured airtightness is often unavailable for an individual dwelling.
In such cases, the simulation depends on the default permeability assigned by
the model. Australian measurements show that residential airtightness varies
substantially with construction period, and that values representative of
Australian dwellings can differ considerably from the defaults originally
used by the engine.

The implementation corrections made in Stage~1 establish the infiltration
path itself. In particular, infiltration is included in addition to
purpose-provided ventilation, the entering air is assigned the correct outdoor
temperature, and leakage is calculated only across the air-exposed envelope.
These corrections are retained unchanged in Stage~3. The Australian refinement
therefore does not alter how infiltration is calculated; it changes the
permeability value supplied to that calculation.

Where a measured blower-door value is available, the measured permeability is
used directly. Where no measurement is available, an age-dependent default is
assigned using Australian residential evidence. The adopted default values are
shown in Table~\ref{tab:AU_leakage}.

\begin{table}[htbp]
\centering
\caption{Australian envelope-permeability defaults used where a measured value
is unavailable. Values are expressed at 50~Pa.}
\label{tab:AU_leakage}
\begin{tabular}{lc}
\toprule
\textbf{Construction period} &
\textbf{$q_{50}$ [m$^3$/(h\,m$^2$)]} \\
\midrule
Before 1900    & 30.0 \\
1901--1920     & 28.0 \\
1921--1945     & 26.0 \\
1946--1960     & 24.0 \\
1961--1975     & 20.0 \\
1976--1990     & 16.0 \\
1991--2005     & 14.0 \\
2006--present  & 11.0 \\
\bottomrule
\end{tabular}
\end{table}

For the case-study building, which was completed in 2006, the adopted value is
treated cautiously because the building lies at the boundary between two
construction periods and no measured airtightness result is available.
The influence of this assumption is therefore examined through sensitivity
testing rather than interpreting a single default value as a measured property
of the dwelling.

The infiltration rate continues to respond to the hourly weather conditions
through the existing infiltration model. The Stage~3 change is therefore
limited to the Australian-specific characterisation of envelope permeability
and does not introduce an independent wind model. Wind-dependent exterior
convection remains the separate Stage~2 mechanism.

% ---------------------------------------------------
\subsubsection{Australian construction-layer representation}
% ---------------------------------------------------

The second refinement concerns the representation of opaque constructions.
As discussed in Stage~2, previous studies have identified limitations in the
simplified treatment of thermal mass within ISO~52016-1, particularly for
heavyweight constructions. That method-level modification is not adopted
directly in this study because it is less relevant to the predominantly
lightweight residential construction targeted here.

The Australian refinement addresses a different issue. Instead of changing
the ISO~52016-1 procedure used to calculate heat transfer and distribute
thermal capacity across the calculation nodes, the physical construction
supplied to that procedure is described using material layers representative
of Australian residential practice.

This distinction separates two levels of modelling. The literature-based
thermal-mass refinement changes \emph{how} a given construction is represented
mathematically within the simplified method. The present refinement changes
\emph{which construction} is supplied to the model. The ISO~52016-1 calculation
procedure itself is retained.

For the initial diagnostic comparison in Stages~1 and~2, the opaque elements
are intentionally assigned zero heat capacity. This removes heat storage and
makes implementation errors and boundary-condition effects easier to identify.
The construction-layer analysis therefore uses a subsequent configuration in
which material-specific thermal properties are restored. This prevents the
effect of construction storage from being mixed with the implementation
diagnosis performed in the earlier stages.

For the Australian configuration, opaque elements are represented through
their constituent material layers rather than only through a generic
construction description. In the case-study external wall, for example, the
brick-veneer assembly is represented according to its actual layer arrangement.
This allows the location of resistance and thermal capacity within the
construction to be represented more realistically while retaining the
ISO~52016-1 calculation framework.

The purpose of this refinement is not to reproduce the detailed construction
model used by EnergyPlus node-for-node. Instead, it tests whether supplying
the simplified model with a construction description that better reflects
Australian residential practice changes its predicted heating and cooling
energy need and its agreement with the reference model.

% ---------------------------------------------------
\subsubsection{Evaluation of the Australian-specific refinements}
% ---------------------------------------------------

The two Australian-specific refinements are evaluated after the preceding
corrections have been established. A common corrected configuration is used
as the reference state, after which the Australian leakage parameterisation
and the Australian construction-layer representation are introduced in
controlled comparisons.

The refinements are examined separately before being combined. This allows a
change caused by the assumed envelope permeability to be distinguished from a
change caused by construction representation. The comparison therefore
includes the configurations shown in Table~\ref{tab:AU_variants}.

\begin{table}[htbp]
\centering
\caption{Configurations used to evaluate the Australian-specific refinements.}
\label{tab:AU_variants}
\begin{tabular}{p{4.0cm}p{4.3cm}p{4.3cm}}
\toprule
\textbf{Configuration} &
\textbf{Envelope leakage} &
\textbf{Construction representation} \\
\midrule

Corrected reference &
Original/default parameterisation &
Reference construction \\

Australian leakage only &
Australian age-dependent permeability &
Reference construction \\

Australian construction only &
Original/default parameterisation &
Australian layer-based construction \\

Combined Australian model &
Australian age-dependent permeability &
Australian layer-based construction \\
\bottomrule
\end{tabular}
\end{table}

All other modelling assumptions are held unchanged within each controlled
comparison. Heating need, cooling need and total energy need are then compared
with the corresponding EnergyPlus reference. The individual comparisons show
the influence of each regional refinement, while the combined configuration
represents the final Australian-adapted version of the model.

The detailed material properties, construction-layer definitions, permeability
sources and infiltration equations are provided in Appendix~C. The Results
section reports the effect of the leakage and construction refinements
separately before presenting the performance of the combined Australian model.








